\section*{Scope and reading guide}
This blueprint is a curated mathematical map to the existing Lean declarations, not a complete transcription of every supporting lemma. The graph records explanatory dependencies; it is not an automatically extracted proof-term dependency graph. The English repository is the canonical formal source. The Japanese companion explains a fixed commit of that source and contains no duplicate proof library.

Read the independent domain first, then the two dimension branches and the four final theorems. Each Lean button resolves to an exact source declaration. These pages are a source browser, not doc-gen API documentation. A green node records a declaration present in the checked build and an editorial correspondence to the prose. It does not certify the accuracy of that prose. Independent human mathematical sign-off has not been recorded.

The domain consists only of isolated hypersurface canonical-root algebras over the complex numbers, not all isolated hypersurfaces. The variable count is the minimum algebra-generator count; the ring dimension is one less. Compiled execution, the JSON text parser/printer, OS output and external hashes remain outside the pure Lean payload theorem.

\section{Objects and common algebra}

\begin{definition}[The actual graded root algebra]
\label{def:root}
\lean{CanonicalRoots.DegreeGroup,CanonicalRoots.IsCanonicalRoot,CanonicalRoots.AmbientRing,CanonicalRoots.RootRing,CanonicalRoots.rootPiece,CanonicalRoots.GradedAlgEquiv}
\leanok
For $p_i\geq2$ with $\sum_i1/p_i<1$, form
\[L=(\mathbb Z^n\oplus\mathbb Zc)/\langle p_ix_i-c\rangle,\quad
T=\mathbb C[X_1,\ldots,X_n]/(\sum_iX_i^{p_i}),\quad\omega=c-\sum_i x_i.\]
A root is an actual $\tau\in L$ with $a\tau=\omega$; its algebra is
$R=\bigoplus_{m\geq0}T_{m\tau}$, realized as a subalgebra of $T$.
The degree group retains torsion. Equivalence means a complex algebra equivalence preserving every graded piece.
\end{definition}

\begin{definition}[The independent classification domain]
\label{def:target}
\lean{CanonicalRoots.RootHypersurfacePresentation,CanonicalRoots.Target}
\leanok
\uses{def:root}
A target for $n\geq3$, $a\geq1$ consists of an admissible signature, an actual root, and a presentation $\mathbb C[z_1,\ldots,z_n]/(f)\simeq R$ as graded algebras. The weights are positive, $f$ is nonzero and homogeneous of positive degree, has no constant or linear terms, and its gradient vanishes only at the origin over all complex points. No enumeration membership or normal-form assumption occurs in this definition.
\end{definition}

\begin{theorem}[Root existence in the integral degree group]
\label{thm:root-arithmetic}
\lean{CanonicalRoots.canonicalRoot_exists_iff}
\leanok
For positive $p_i$, let $N=\operatorname{lcm}(p_i)$ and $k=N-\sum_iN/p_i$. An actual canonical root exists if and only if $a\mid k$ and $\gcd(a,N)=1$.
\end{theorem}
\begin{proof}
\leanok
\uses{def:root}
Evaluate the root equation in the degree-group normal form. The residue equations force coprimality, and the numerator degree forces divisibility. Conversely, Bezout coefficients construct the root in the original quotient group.
\end{proof}

\begin{definition}[A finite free basis]
\label{def:basis}
\lean{CanonicalRoots.rootFiniteFreeBasis}
\leanok
\uses{thm:root-arithmetic}
For an admissible root and positive integers $N,u$ satisfying $N\tau=uc$, the root algebra has an explicit finite basis over the subalgebra generated by the specified coordinate powers. The linked declaration includes the precise box index and scalar action.
\end{definition}

\begin{theorem}[The Hilbert series from actual components]
\label{thm:hilbert}
\lean{CanonicalRoots.canonicalRootHilbertSeries_eq_rational}
\leanok
For an admissible canonical root, its Hilbert series equals the rational expression obtained from the finite box of monomials and the polynomial-power denominator.
\end{theorem}
\begin{proof}
\leanok
\uses{def:basis}
Regroup the homogeneous basis by box residues and nonnegative powers; count degrees to identify the power-series coefficients and then the rational expression.
\end{proof}

\begin{theorem}[Intrinsic numerical identities]
\label{thm:numerics}
\lean{CanonicalRoots.Target.relationDegree_eq_add_sum_weights,CanonicalRoots.Target.weights_gcd_eq_one}
\leanok
Every target has primitive weights and relation degree $h=a+\sum_iw_i$. These are derived properties of the actual root presentation.
\end{theorem}
\begin{proof}
\leanok
\uses{def:target,thm:hilbert}
Compare the root Hilbert expression with the hypersurface Hilbert expression. The reciprocal identity gives the degree difference; eventual nonvanishing of graded pieces forces the weight gcd to be one.
\end{proof}

\begin{theorem}[Exactly the minimum number of generators]
\label{thm:minimum}
\lean{CanonicalRoots.RootHypersurfacePresentation.root_minimum_generators,CanonicalRoots.RootHypersurfacePresentation.exists_minimal_surjection}
\leanok
A root presentation in $n$ variables yields a surjection from an $n$-variable polynomial algebra. Any surjection from an $m$-variable polynomial algebra onto that root algebra has $n\leq m$, even without requiring the alternative map to preserve degrees.
\end{theorem}
\begin{proof}
\leanok
\uses{def:target}
The missing constant and linear terms give an $n$-dimensional cotangent space at the origin. A surjective presentation bounds that dimension by $m$; transfer through the actual algebra equivalence.
\end{proof}

\begin{theorem}[The canonical-module shift]
\label{thm:canonical}
\lean{CanonicalRoots.Target.canonical_parameter}
\leanok
Every target satisfies the canonical-parameter contract: the actual principal $\operatorname{Ext}^1$ module, with its descended quotient action, has graded pieces corresponding to $R_{m+a}$ for every integer $m$.
\end{theorem}
\begin{proof}
\leanok
\uses{thm:numerics}
Use the principal free resolution and its graded Hom construction to identify the actual Ext module. Substitute $h=a+\sum w_i$ into its degree formula. The parameter is not inserted into the definition of the canonical grading.
\end{proof}

\section{The ternary branch}

\begin{theorem}[The complete ternary arithmetic search]
\label{thm:ternary-enum}
\lean{CanonicalRoots.enumerateTernaryCandidates_iff}
\leanok
For $a\geq1$, membership in the finite ternary candidate list is equivalent to the independent arithmetic predicate: one of the five types, exponents at least two, primitive unreduced weights, and defect $a$. The search box has exponent bound $a+6$.
\end{theorem}
\begin{proof}
\leanok
\uses{def:root}
Prove the uniform exponent bound for each type, then identify membership in the finite box and its filter. This arithmetic statement alone is not semantic completeness.
\end{proof}

\begin{theorem}[From an arbitrary ternary target to five types]
\label{thm:ternary-reduction}
\lean{CanonicalRoots.Target.ternary_arithmetic_candidate}
\leanok
Every target with three variables admits an arithmetic ternary candidate with the same relation degree and weights up to permutation.
\end{theorem}
\begin{proof}
\leanok
\uses{thm:numerics}
The implementation follows actual Hilbert poles, signature recovery, support analysis, exclusion of branched cases, and primitive normalization. Apply these results to the target, then use primitive weights and the degree identity to meet the arithmetic predicate.
\end{proof}

\begin{theorem}[Actual ternary realizations and completeness]
\label{thm:ternary-model}
\lean{CanonicalRoots.ternary_candidate_realization,CanonicalRoots.Target.ternary_enumerated_model}
\leanok
Every arithmetic ternary candidate realizes an actual canonical-root target. Conversely every ternary target is graded-isomorphic to a model from the enumerated candidates.
\end{theorem}
\begin{proof}
\leanok
\uses{thm:ternary-reduction,thm:ternary-enum}
Construct the five monomial quotient maps and prove injectivity and surjectivity through their basis and reduction arguments. Combine these realizations with the target-to-candidate theorem and the root classification by the recovered data. Hilbert-series equality alone is not used as a general isomorphism principle.
\end{proof}

\begin{theorem}[The ternary key classifies actual models]
\label{thm:key}
\lean{CanonicalRoots.candidateKey_eq_iff_gradedEquiv}
\leanok
For two arithmetic candidates at the same positive parameter, their sorted weight-and-degree keys agree if and only if their presented graded complex algebras are isomorphic.
\end{theorem}
\begin{proof}
\leanok
\uses{thm:ternary-model,thm:minimum}
A matching key gives a weight permutation, and the ternary root classification gives an actual graded equivalence. Conversely a graded equivalence preserves the minimal generator weights; the common parameter determines the relation degree.
\end{proof}

\section{The higher-dimensional branch}

\begin{definition}[The character action on the actual quotient]
\label{def:action}
\lean{CanonicalRoots.ambientCharacterAction,CanonicalRoots.ambientCharacterFixed}
\leanok
\uses{def:root}
Characters of $L/\mathbb Z\tau$ act diagonally on the ambient quotient $T$. The fixed subalgebra consists of elements fixed by every such automorphism. The relation is preserved before the action is descended to the quotient.
\end{definition}

\begin{theorem}[The fixed-ring identification]
\label{thm:fixed}
\lean{CanonicalRoots.rootSubalgebra_eq_characterFixed,CanonicalRoots.rootRingEquivCharacterFixed}
\leanok
For an admissible canonical root at positive parameter, the actual root subalgebra equals the character-fixed subalgebra: $R=T^G$. In particular there is a complex algebra equivalence between them.
\end{theorem}
\begin{proof}
\leanok
\uses{def:action,thm:root-arithmetic}
Every root piece is fixed. In the other direction, decompose a fixed element into homogeneous components. Character separation forces each nonzero component into the integer root line; positivity makes the negative pieces vanish.
\end{proof}

\begin{theorem}[Higher-dimensional coprimality]
\label{thm:coprime}
\lean{CanonicalRoots.Target.pairwise_coprime}
\leanok
For every target in at least four variables, the signature entries are pairwise coprime.
\end{theorem}
\begin{proof}
\leanok
\uses{def:target,thm:fixed}
A common divisor produces a nonzero stratum with a stabilizer quotient. The formal invariant chart and cotangent-space lower bound contradict the regularity away from the vertex of an isolated presentation. The stabilizer, rather than the entire group, is used in the local chart.
\end{proof}

\begin{theorem}[The higher-dimensional criterion]
\label{thm:higher-criterion}
\lean{CanonicalRoots.higher_root_hypersurface_iff,CanonicalRoots.RootHypersurfacePresentation.higher_root_eq_top}
\leanok
For an admissible actual root at positive parameter in at least four variables, an isolated hypersurface presentation exists exactly when the signature is pairwise coprime and
$P-\sum_iP/p_i=a$, where $P=\prod_ip_i$. In that case the root subalgebra is all of $T$.
\end{theorem}
\begin{proof}
\leanok
\uses{thm:coprime,thm:minimum}
After coprimality, pure-power generation and support counting force root scale one. Applying the product-degree map to the root equation gives the defect. Conversely the Fermat presentation realizes the arithmetic data.
\end{proof}

\begin{theorem}[A complete integer recursion]
\label{thm:higher-enum}
\lean{CanonicalRoots.enumerateHigherCandidates_iff}
\leanok
For $a>0$, the higher candidate list is exactly the strictly increasing lists of length $n$, with entries at least two, pairwise coprime, and defect $a$. The recursion decreases the number of remaining entries.
\end{theorem}
\begin{proof}
\leanok
\uses{def:root}
Keep prefix product $P$ and integer defect $A$. With $r$ entries left the next entry is bounded by $\lfloor(r+a)P/A\rfloor$; update $(P,A)$ to $(Pq,Aq-P)$. The last entry is forced by $Aq=a+P$. Induction proves that every valid prefix stays in the search.
\end{proof}

\begin{theorem}[The higher-dimensional semantic classification]
\label{thm:higher-model}
\lean{CanonicalRoots.higher_candidates_sound,CanonicalRoots.Target.higher_candidates_complete_unique,CanonicalRoots.higher_candidates_pairwise_nonisomorphic}
\leanok
The higher candidate models are actual targets. Every higher-dimensional target is isomorphic to exactly one output candidate, and different candidate positions are not graded-isomorphic.
\end{theorem}
\begin{proof}
\leanok
\uses{thm:higher-criterion,thm:higher-enum}
Combine the independent hypersurface criterion with the complete integer recursion, signature reindexing, and the intrinsic graded invariants of the Fermat models.
\end{proof}

\section{The classifier and final four theorems}

\begin{definition}[The public total classifier]
\label{def:enumerate}
\lean{CanonicalRoots.Input,CanonicalRoots.enumerate}
\leanok
\uses{thm:ternary-enum,thm:higher-enum}
An input contains $n,a$ and proofs only of $3\leq n$ and $1\leq a$. The total function returns a finite list of exact equation data, choosing the ternary or higher branch according to $n$.
\end{definition}

\begin{theorem}[Transport to actual output equations]
\label{thm:output}
\lean{CanonicalRoots.Target.output_complete_unique,CanonicalRoots.output_pairwise_nonisomorphic}
\leanok
Every target occurs at exactly one position in the equation-data list up to graded algebra equivalence; distinct positions are nonisomorphic.
\end{theorem}
\begin{proof}
\leanok
\uses{thm:key,thm:higher-model,def:enumerate}
Split on the dimension and use the semantic classifications. Transport the candidate models through the synchronized variable permutation used to sort the output weights.
\end{proof}

\begin{theorem}[The exact root witness in every row]
\label{thm:witness}
\lean{CanonicalRoots.output_root_witness}
\leanok
Every output row has an actual target whose polynomial, weights, relation degree, ordered signature and canonical root agree with the displayed data.
\end{theorem}
\begin{proof}
\leanok
\uses{thm:output,thm:root-arithmetic}
Refine output realization with the original signature and transfer the root normal-form witness. The proof also accounts for the permutation of weights and exponent columns.
\end{proof}

\begin{definition}[The public soundness contract]
\label{def:contract}
\lean{CanonicalRoots.EquationRealizes}
\leanok
\uses{def:target,thm:canonical,thm:minimum}
Equation realization includes dimension and parameter, sorted positive primitive weights, nonnegative exponent data, an actual target with matching polynomial and root witness, the canonical Ext shift, and the exact minimum number of algebra generators. This predicate is defined independently of list membership.
\end{definition}

\begin{theorem}[Final theorem: soundness]
\label{thm:sound}
\lean{CanonicalRoots.enumerate_sound}
\leanok
For every legal input and every equation in its output, the equation-realization contract holds.
\end{theorem}
\begin{proof}
\leanok
\uses{thm:witness,def:contract}
Convert list membership into a position, obtain the exact root witness, and combine the output well-formedness lemmas with the canonical parameter and minimum-generator results.
\end{proof}

\begin{theorem}[Final theorem: unique completeness]
\label{thm:complete}
\lean{CanonicalRoots.enumerate_complete}
\leanok
For every legal input and every independently defined target at that input, there exists exactly one output position whose presented ring is graded-isomorphic to the actual root ring.
\end{theorem}
\begin{proof}
\leanok
\uses{def:target,thm:output}
Apply the semantic completeness and uniqueness theorem for the actual equation-data list.
\end{proof}

\begin{theorem}[Final theorem: pairwise nonisomorphism]
\label{thm:noniso}
\lean{CanonicalRoots.enumerate_pairwise_nonisomorphic}
\leanok
Two distinct positions in the output list cannot be isomorphic as graded complex algebras.
\end{theorem}
\begin{proof}
\leanok
\uses{thm:output}
Use the actual output nonisomorphism theorem, not merely inequality of executable keys.
\end{proof}

\begin{theorem}[Final theorem: pure payload correctness]
\label{thm:payload}
\lean{CanonicalRoots.cli_payload_correct}
\leanok
Decoding the pure normal classification payload gives exactly $\operatorname{some}(n,a,\operatorname{enumerate}(\mathrm{input}))$. This concerns the pure JSON value and decoder, not execution of the OS output channel.
\end{theorem}
\begin{proof}
\leanok
\uses{def:enumerate}
Apply the generic classification-payload encoding/decoding round trip to the complete equation list.
\end{proof}

\begin{theorem}[Checking a whole output list]
\label{thm:checker}
\lean{CanonicalRoots.checkEquationList_correct}
\leanok
The equation-list checker returns true exactly when the supplied list equals the classifier result, including order and multiplicities.
\end{theorem}
\begin{proof}
\leanok
\uses{def:enumerate,thm:payload}
Unfold the Boolean equality check and use decidable equality of the full equation-data lists. Concrete whole-list certificates are supplied in the separate OutputCertificates module.
\end{proof}

\section*{Verification and sources}
Run the repository verification script and inspect the final theorem signatures as described in docs/VERIFICATION.md. The preserved research manuscript and historical self-audit are in the English repository; they are proof candidates and provenance, not axioms. The build uses Patrick Massot's \href{https://github.com/PatrickMassot/leanblueprint}{Lean Blueprint} package with plasTeX and Graphviz.
